\chapter{الهندسة التحليلية}\label{ch:g10:coordgeom}

كانت فكرة ديكارت الكبرى أن يصف النقط بأزواج من الأعداد، فيحوّل
مسائل الهندسة إلى حسابات. وبصيغتين اثنتين فقط --- صيغة \omterm{prop:g10:coordgeom:midpoint}{المنتصف}
وصيغة المسافة --- يمكن البرهان على أن مثلثًا متساوي الساقين،
أو أن رباعيًا متوازي أضلاع، أو أن ثلاث نقط على استقامة واحدة،
دون رسم أي مستقيم مساعد.

\section{الإحداثيات في المستوي}

\begin{definition}[نظام الإحداثيات]\label{def:g10:coordgeom:system}
يتألف \emph{نظام الإحداثيات}\index{نظام إحداثيات} في المستوي
من مبدأ $O$ ومحورين مدرّجين يمران به: المحور الأفقي
$x$ والمحور الشاقولي $y$. وعندئذٍ يكون لكل نقطة $M$ زوج وحيد من
\emph{الإحداثيات}\index{إحداثيات} $(x, y)$: هما
\emph{فاصلتها}\index{فاصلة} $x$ و\emph{ترتيبها}\index{ترتيب} $y$.
ويكون النظام \emph{متعامدًا ومتجانسًا}\index{نظام متعامد ومتجانس} إذا كان المحوران
متعامدين ويحملان وحدة الطول نفسها. وكل الأنظمة في هذا
الفصل متعامدة ومتجانسة.
\end{definition}

\begin{omfigure}
\begin{tikzpicture}[scale=0.85]
  \draw[thin, gray!40] (-2.4,-1.4) grid (4.4,3.4);
  \draw[-Stealth] (-2.6,0) -- (4.6,0) node[below, font=\small] {$x$};
  \draw[-Stealth] (0,-1.6) -- (0,3.6) node[left, font=\small] {$y$};
  \node[below left, font=\small] at (0,0) {$O$};
  \fill[omDef] (3,2) circle (2pt);
  \node[above right, font=\small, omDef] at (3,2) {$M(3, 2)$};
  \draw[dashed, omDef] (3,0) -- (3,2) -- (0,2);
  \fill[omThm] (-2,-1) circle (2pt);
  \node[below left, font=\small, omThm] at (-2,-1) {$N(-2, -1)$};
  \foreach \x in {1,2,3,4} \node[below, font=\footnotesize] at (\x,0) {\x};
  \foreach \y in {1,2,3} \node[left, font=\footnotesize] at (0,\y) {\y};
\end{tikzpicture}

{\small تُحدَّد كل نقطة من المستوي بعددين: \omterm{def:g10:coordgeom:system}{الفاصلة} أولًا
(أفقيًا)، ثم \omterm{def:g10:coordgeom:system}{الترتيب} (شاقوليًا).}
\end{omfigure}

\section{منتصف قطعة}

\begin{proposition}[صيغة المنتصف]\label{prop:g10:coordgeom:midpoint}
لتكن $A(x_A, y_A)$ و $B(x_B, y_B)$. \omterm{def:g10:coordgeom:system}{إحداثيات} \emph{منتصف}\index{منتصف}
القطعة $[AB]$، وليكن $I$، هي
\[
I\left(\frac{x_A + x_B}{2},\ \frac{y_A + y_B}{2}\right).
\]
\end{proposition}

\begin{proof}
لننظر في \omterm{def:g10:coordgeom:system}{الإحداثيات} الأفقية. عند الانتقال من $A$ إلى $B$، تتغير \omterm{def:g10:coordgeom:system}{الفاصلة}
بمقدار $x_B - x_A$؛ إذن \omterm{def:g10:coordgeom:system}{فاصلة} النقطة الواقعة في \omterm{prop:g10:coordgeom:midpoint}{منتصف} الطريق هي
\[
x_A + \frac{x_B - x_A}{2} = \frac{2x_A + x_B - x_A}{2}
= \frac{x_A + x_B}{2},
\]
أي متوسط الفاصلتين. وينطبق الحساب نفسه على
التراتيب.
\end{proof}

\begin{example}\label{ex:g10:coordgeom:midpoint}
من أجل $A(-1, 4)$ و $B(5, -2)$، يكون \omterm{prop:g10:coordgeom:midpoint}{منتصف} $[AB]$
\[
I\left(\frac{-1 + 5}{2},\ \frac{4 + (-2)}{2}\right) = I(2, 1).
\]
وتعمل الصيغة في الاتجاه المعاكس أيضًا: إذا كانت $A(-1, 4)$ \omterm{prop:g10:coordgeom:midpoint}{والمنتصف}
$I(2, 1)$، فإن $B$ تحقق $\frac{-1 + x_B}{2} = 2$ و
$\frac{4 + y_B}{2} = 1$، إذن $x_B = 5$ و $y_B = -2$: أي أن النقطة $B$ هي
\emph{نظيرة} $A$ بالنسبة إلى $I$.
\end{example}

\section{المسافة بين نقطتين}

\begin{theorem}[صيغة المسافة]\label{thm:g10:coordgeom:distance}
في \omterm{def:g10:coordgeom:system}{نظام متعامد ومتجانس}، المسافة بين $A(x_A, y_A)$ و
$B(x_B, y_B)$ هي
\[
AB = \sqrt{(x_B - x_A)^2 + (y_B - y_A)^2}.
\]
\end{theorem}

\begin{proof}
لتكن $C$ النقطة $(x_B, y_A)$: لها ارتفاع $A$ نفسه \omterm{def:g10:coordgeom:system}{وفاصلة}
$B$ نفسها، إذن للمثلث $ACB$ زاوية قائمة في $C$. وضلعا
القائمة أفقي وشاقولي، وطولاهما
$AC = \abs{x_B - x_A}$ و $CB = \abs{y_B - y_A}$. وحسب مبرهنة
فيثاغورس،
\[
AB^2 = AC^2 + CB^2 = (x_B - x_A)^2 + (y_B - y_A)^2,
\]
وبأخذ الجذر التربيعي (والطرفان غير سالبين) نحصل على الصيغة.
\end{proof}

\begin{omfigure}
\begin{tikzpicture}[scale=0.85]
  \draw[thin, gray!40] (-0.4,-0.4) grid (5.4,3.4);
  \draw[-Stealth] (-0.6,0) -- (5.6,0) node[below, font=\small] {$x$};
  \draw[-Stealth] (0,-0.6) -- (0,3.6) node[left, font=\small] {$y$};
  \coordinate (A) at (1,1);
  \coordinate (B) at (5,3);
  \coordinate (C) at (5,1);
  \draw[omDef, thick] (A) -- (B);
  \draw[omThm, thick] (A) -- (C) node[midway, below, font=\small]
    {$\abs{x_B - x_A}$};
  \draw[omThm, thick] (C) -- (B) node[midway, right, font=\small]
    {$\abs{y_B - y_A}$};
  \draw[thin] (C) ++(-0.25,0) -- ++(0,0.25) -- ++(0.25,0);
  \fill (A) circle (2pt) node[above left, font=\small] {$A$};
  \fill (B) circle (2pt) node[above, font=\small] {$B$};
  \fill (C) circle (2pt) node[below right, font=\small] {$C$};
\end{tikzpicture}

{\small صيغة المسافة هي مبرهنة فيثاغورس مطبقة على
المثلث القائم $ACB$ المبني على ضلع أفقي وضلع شاقولي.}
\end{omfigure}

\begin{example}\label{ex:g10:coordgeom:distance}
من أجل $A(1, 1)$ و $B(5, 3)$:
\[
AB = \sqrt{(5-1)^2 + (3-1)^2} = \sqrt{16 + 4} = \sqrt{20} = 2\sqrt 5 .
\]
وتُترك المسافات عادةً على صورتها المضبوطة (بجذر تربيعي)؛ ولا نقرّب إلا
في النهاية إذا لزم جواب عشري.
\end{example}

\begin{remark}
لا يهم \omterm{def:g10:coordgeom:system}{ترتيب} النقطتين: $(x_A - x_B)^2 =
(x_B - x_A)^2$. لكن لا تنسَ التربيع! فالمسافة
\emph{ليست} $\abs{x_B - x_A} + \abs{y_B - y_A}$.
\end{remark}

\section{استعمال الصيغتين في الهندسة}

\begin{method}[طبيعة مثلث]\label{met:g10:coordgeom:triangle}
إذا أُعطيت ثلاث نقط $A$ و $B$ و $C$ بإحداثياتها:
\begin{enumerate}
  \item احسب مربعات الأطوال الثلاثة $AB^2$ و $AC^2$ و $BC^2$ بصيغة
        المسافة (فالاحتفاظ بالمربعات يجنّبك الجذور)؛
  \item إذا تساوى مربعا طولين $\Rightarrow$ فالمثلث
        \emph{متساوي الساقين}؛
  \item وإذا كان أكبر مربعات الأطوال مساويًا لمجموع الآخرين، فالمثلث
        \emph{قائم الزاوية} في الرأس المقابل للضلع
        الأطول، حسب عكس مبرهنة فيثاغورس.
\end{enumerate}
\end{method}

\begin{example}\label{ex:g10:coordgeom:triangle}
لتكن $A(0, 1)$ و $B(4, 3)$ و $C(2, -3)$. عندئذٍ
\begin{align*}
AB^2 &= (4-0)^2 + (3-1)^2 = 16 + 4 = 20, \\
AC^2 &= (2-0)^2 + (-3-1)^2 = 4 + 16 = 20, \\
BC^2 &= (2-4)^2 + (-3-3)^2 = 4 + 36 = 40 .
\end{align*}
أولًا، $AB^2 = AC^2$، إذن $AB = AC$: فالمثلث متساوي الساقين في $A$.
وعلاوة على ذلك $AB^2 + AC^2 = 40 = BC^2$، إذن يكون المثلث كذلك قائم الزاوية في $A$
حسب عكس مبرهنة فيثاغورس.
\end{example}

\begin{method}[التعرّف على متوازي الأضلاع]\label{met:g10:coordgeom:parallelogram}
يكون الرباعي $ABCD$ متوازي أضلاع بالضبط عندما يكون لقطريه
$[AC]$ و $[BD]$ \omterm{prop:g10:coordgeom:midpoint}{المنتصف} نفسه. إذن: احسب المنتصفين بصيغة
\omterm{prop:g10:coordgeom:midpoint}{المنتصف} وقارن بينهما.
\end{method}

\begin{example}\label{ex:g10:coordgeom:parallelogram}
لتكن $A(-1, 0)$ و $B(2, 2)$ و $C(5, 1)$ و $D(2, -1)$. \omterm{prop:g10:coordgeom:midpoint}{منتصف}
$[AC]$ هو $\left(\frac{-1+5}{2}, \frac{0+1}{2}\right) = (2, \frac12)$؛ \omterm{prop:g10:coordgeom:midpoint}{ومنتصف}
$[BD]$ هو $\left(\frac{2+2}{2}, \frac{2-1}{2}\right) =
(2, \frac12)$. وهما متطابقان، إذن $ABCD$ متوازي أضلاع.
\end{example}

\begin{omfigure}
\begin{tikzpicture}[scale=0.75]
  \draw[thin, gray!40] (-1.4,-1.4) grid (5.4,2.4);
  \draw[-Stealth] (-1.6,0) -- (5.6,0);
  \draw[-Stealth] (0,-1.6) -- (0,2.6);
  \coordinate (A) at (-1,0);
  \coordinate (B) at (2,2);
  \coordinate (C) at (5,1);
  \coordinate (D) at (2,-1);
  \draw[omDef, thick] (A) -- (B) -- (C) -- (D) -- cycle;
  \draw[omThm, dashed] (A) -- (C);
  \draw[omThm, dashed] (B) -- (D);
  \fill (A) circle (2pt) node[left, font=\small] {$A$};
  \fill (B) circle (2pt) node[above, font=\small] {$B$};
  \fill (C) circle (2pt) node[right, font=\small] {$C$};
  \fill (D) circle (2pt) node[below, font=\small] {$D$};
  \fill[omThm] (2,0.5) circle (2pt) node[above right, font=\small] {$I$};
\end{tikzpicture}

{\small الرباعي $ABCD$ متوازي أضلاع لأن لقطريه \omterm{prop:g10:coordgeom:midpoint}{المنتصف}
نفسه $I(2, \frac12)$.}
\end{omfigure}

\section{تمارين}

\begin{exercise}[$\star$]\label{exo:g10:coordgeom:1}
لتكن $A(2, 5)$ و $B(-4, 1)$. احسب \omterm{def:g10:coordgeom:system}{إحداثيات} \omterm{prop:g10:coordgeom:midpoint}{منتصف}
$[AB]$ والمسافة $AB$.
\end{exercise}

\begin{exercise}[$\star$]\label{exo:g10:coordgeom:2}
لتكن $A(3, -2)$ و $I(1, 2)$. أوجد \omterm{def:g10:coordgeom:system}{إحداثيات} النقطة $B$ التي
تجعل $I$ \omterm{prop:g10:coordgeom:midpoint}{منتصف} $[AB]$.
\end{exercise}

\begin{exercise}[$\star$]\label{exo:g10:coordgeom:3}
أي النقط $P(4, 1)$ و $Q(-3, 2)$ و $R(0, -5)$ أقرب إلى
المبدأ $O(0,0)$؟ أجب بمسافات مضبوطة.
\end{exercise}

\begin{exercise}[$\star$]\label{exo:g10:coordgeom:4}
علّم النقط $A(1, 2)$ و $B(5, 4)$ و $C(7, 0)$ على ورق مرسوم، ثم بيّن
بالحساب أن المثلث $ABC$ متساوي الساقين. وهل هو قائم الزاوية؟
\end{exercise}

\begin{exercise}[$\star\star$]\label{exo:g10:coordgeom:5}
لتكن $A(-2, 1)$ و $B(1, 3)$ و $C(4, 1)$ و $D(1, -1)$.
\begin{enumerate}
  \item بيّن أن $ABCD$ متوازي أضلاع.
  \item احسب $AB$ و $BC$. هل $ABCD$ معيّن (كل أضلاعه متساوية)؟
  \item احسب $AC$ و $BD$. هل $ABCD$ مستطيل (قطراه متساويان)؟
\end{enumerate}
\end{exercise}

\begin{exercise}[$\star\star$]\label{exo:g10:coordgeom:6}
الدائرة التي مركزها $\Omega(2, 1)$ ونصف قطرها $5$ هي مجموعة كل
النقط التي تبعد $5$ عن $\Omega$. أي النقط $A(5, 5)$ و
$B(-1, -3)$ و $C(6, 2)$ تقع على هذه الدائرة؟ وأيها داخلها؟ وأيها خارجها؟
\end{exercise}

\begin{exercise}[$\star\star$]\label{exo:g10:coordgeom:7}
لتكن $A(1, 1)$ و $B(7, 5)$. أوجد كل نقط $M(x, 0)$ من المحور $x$
المتساوية البعد عن $A$ و $B$. (اكتب $MA^2 = MB^2$ وحل بالنسبة إلى
$x$.)
\end{exercise}

\begin{exercise}[$\star\star$]\label{exo:g10:coordgeom:8}
لتكن $A(0, 3)$ و $B(4, 1)$.
\begin{enumerate}
  \item احسب \omterm{def:g10:coordgeom:system}{إحداثيات} \omterm{prop:g10:coordgeom:midpoint}{منتصف} $[AB]$، وليكن $I$.
  \item بيّن أن $M(2, 2) = I$، ثم تحقق بحساب $MA$ و $MB$ من أن
        $M$ متساوية البعد عن $A$ و $B$.
  \item أوجد نقطة ثانية، على المحور $y$، متساوية البعد عن $A$ و
        $B$.
\end{enumerate}
\end{exercise}

\begin{exercise}[$\star\star$]\label{exo:g10:coordgeom:9}
أُعطيت النقط $A(-1, -1)$ و $B(3, 1)$ و $C(11, 5)$. احسب $AB$
و $BC$ و $AC$، واستنتج أن $A$ و $B$ و $C$ على استقامة واحدة. (تكون ثلاث نقط
على استقامة واحدة بالضبط عندما تكون أكبر المسافات الثلاث مساويةً لمجموع
الأخريين.)
\end{exercise}

\begin{exercise}[$\star\star\star$]\label{exo:g10:coordgeom:10}
لتكن $A(a, 0)$ و $B(0, b)$ حيث $a, b > 0$، وليكن $I$ \omterm{prop:g10:coordgeom:midpoint}{منتصف}
$[AB]$. بيّن بالحساب أن $OI = IA = IB$، حيث $O$ هو
المبدأ. (وهذا يبرهن على مبرهنة كلاسيكية: في مثلث قائم،
يكون \omterm{prop:g10:coordgeom:midpoint}{منتصف} الوتر متساوي البعد عن الرؤوس الثلاثة.)
\end{exercise}

\section{مسألة: جسر ديكارت --- معادلات الدوائر وكيف
تعرف أين أنت}

\begin{problem}[{مسألة نهاية الأسبوع --- الدائرة تصير
معادلة، والمبرهنات القديمة تصير حسابات، وثلاث منارات
تحدّد نقطة: التموضع بالمسافات جبريًا}]
\label{pb:g10:coordgeom:1}
سنة 1637 بنى رينيه ديكارت جسرًا بين قارتين:
فصار كل \omterm{def:g10:functions:graph}{منحنى} هندسي \emph{\omterm{def:g10:algebra:equation}{معادلة}}، وصار كل
سؤال هندسي حسابًا. تعبر هذه المسألة الجسر
في الاتجاهين --- فتشتق \omterm{def:g10:algebra:equation}{معادلة} الدائرة من
صيغة المسافة (\cref{thm:g10:coordgeom:distance})، وتعيد برهان
المبرهنات الكلاسيكية في ثلاثة أسطر من الجبر، وتنتهي حيث
يحمل الجسر أكبر حركة اليوم: حساب موضع
من إشارات المسافة، وهو قلب التموضع العالمي مستويًا.

\medskip
\noindent\textbf{الجزء الأول --- الدائرة تنال \omterm{def:g10:algebra:equation}{معادلة}.}
\begin{enumerate}
  \item تقع النقطة $M(x, y)$ على الدائرة التي مركزها
        $\Omega(2, 1)$ ونصف قطرها $5$ بالضبط عندما
        $\Omega M = 5$. ربّع هذا الشرط لتحصل على
        \omterm{def:g10:algebra:equation}{معادلة} الدائرة:
        \[
        (x - 2)^2 + (y - 1)^2 = 25 .
        \]
  \item اختبر الانتماء: أي النقط $A(5, 5)$ و $B(6, 4)$ و
        $D(4, 4)$ تقع على الدائرة، وأيها داخلها، وأيها خارجها؟
  \item دائرة متنكرة: أتمّ المربعات
        (\cref{pb:g10:algebra:1}) في
        \[
        x^2 + y^2 - 6x + 4y - 12 = 0
        \]
        وأعطِ مركزها ونصف قطرها.
  \item افعل الشيء نفسه مع $x^2 + y^2 - 6x + 4y + 14 = 0$. ماذا
        تكشف الصورة المكتملة؟ وبيّن المحك: متى
        تصف $(x - h)^2 + (y - k)^2 = c$ دائرة، أم
        نقطة وحيدة، أم لا شيء البتة؟
  \item قاطع دائرة السؤال 3 بالمستقيم الأفقي
        $y = 3$: عوّض وحل. كم نقطة تحصل عليها، وما
        الاسم الهندسي لمستقيم كهذا؟
\end{enumerate}

\medskip
\noindent\textbf{الجزء الثاني --- مبرهنات قديمة في ثلاثة أسطر.}
\begin{enumerate}[resume]
  \item مبرهنة الدائرة في الكتاب السابق (زاوية قائمة في كل نصف دائرة)،
        معادًا برهانها جبريًا: ليكن $A(-r, 0)$ و $B(r, 0)$
        طرفي قطر ولتكن $M(x, y)$ نقطة كيفية من
        الدائرة $x^2 + y^2 = r^2$. احسب $MA^2 + MB^2$
        واختم بعكس مبرهنة فيثاغورس أن الزاوية
        $\widehat{AMB}$ قائمة.
  \item مبرهنة المتوسط: من أجل $A(-a, 0)$ و $B(a, 0)$ (فيكون
        \omterm{prop:g10:coordgeom:midpoint}{منتصف} $[AB]$ هو المبدأ $I$)، بيّن أنه من أجل
        كل نقطة $M(x, y)$:
        \[
        MA^2 + MB^2 = 2\,MI^2 + \frac{AB^2}{2} .
        \]
  \item استنتج المحل الهندسي للنقط $M$ التي تحقق
        $MA^2 + MB^2 = 26$ عندما $AB = 6$: أي \omterm{def:g10:functions:graph}{منحنى}، وأي
        مركز، وأي نصف قطر؟
  \item محل هندسي آخر: $A(0,0)$ و $B(3,0)$؛ أوجد كل النقط
        التي تحقق $MA = 2\,MB$. (ربّع، وانشر، وأتمّ
        المربعات: تظهر دائرة شهيرة --- عرفها أبولونيوس
        دون \omterm{def:g10:coordgeom:system}{إحداثيات}.)
  \item في جملة أو جملتين: ماذا يغيّر جسر ديكارت في
        \emph{كيفية} برهان المبرهنات؟ قارن
        السؤال 6 ببرهان المستطيل في
        مبرهنة الدائرة في الكتاب السابق (زاوية قائمة في كل نصف دائرة).
\end{enumerate}

\medskip
\noindent\textbf{الجزء الثالث --- ثلاث منارات تجدك.}
موضعك $M(x, y)$ مجهول. تقيس المنارة $P(0, 0)$
بعدك فتجده $5$؛ وتقيسه المنارة $Q(6, 0)$ فتجده $5$ أيضًا.
\begin{enumerate}[resume]
  \item اكتب معادلتي الدائرتين، واطرح إحداهما من الأخرى، وراقب
        المربعات تتلاشى: أي \omterm{def:g10:algebra:equation}{معادلة} بسيطة تبقى؟
        اجمعها مع إحدى الدائرتين لتجد الموضعين
        \emph{المرشحين}.
  \item تقيس منارة ثالثة $R(0, 8)$ بعدك فتجده $5$.
        احسب بعدها عن كل مرشح وقرّر أين
        أنت.
  \item فسّر لماذا يعطي طرح معادلتي دائرتين
        \emph{دائمًا} \omterm{def:g10:algebra:equation}{معادلة} مستقيم (أي الحدود
        تتلاشى؟)، وما ذلك المستقيم هندسيًا عندما
        تتقاطع الدائرتان في نقطتين.
  \item يعمل التموضع بالأقمار الصناعية فعليًا في الفضاء وبساعات
        غير مضبوطة: فيعطي كل قمر (عبر زمن انتقال
        الإشارة) مسافةً واحدة، أي كرة واحدة. كم مجهولًا
        لدى المستقبِل (الموضع \emph{وخطأ ساعته} هو)،
        ولماذا يصغي التموضع العالمي إذن إلى \emph{أربعة}
        أقمار على الأقل؟
  \item الدقة: لنفترض أن بعد المنارة $P$ هو في الحقيقة
        $5.1$ بدل $5$ (مع بقاء كل شيء آخر). أعد
        طرح السؤال 11 لتجد $x$ الجديد، وقدّر
        كم انتقل التقدير. (قارن بأفكار ميزانية
        الخطأ في \cref{pb:g10:numbers:1}.)
\end{enumerate}

\medskip
\noindent\textbf{الجزء الرابع --- عدة المساح.}
\begin{enumerate}[resume]
  \item أقصر طريق إلى الماء: مخيّم في $A(1, 5)$، وحظيرة
        في $B(7, 3)$، ونهر مستقيم على المحور $y = 0$.
        للانتقال من $A$ إلى النهر ثم إلى $B$ بأقل
        مشي: اعكس $B$ بالنسبة إلى النهر إلى $B'$،
        وأوجد أين تقطع القطعة $[AB']$ النهر.
        أعطِ نقطة القطع والطول الكلي الأصغري.
        (والفكرة هي تمرين ارتداد البلياردو في الكتاب السابق، وقد صارت الآن
        قابلة للحساب تمامًا.)
  \item أكّد الأصغرية على منافس: احسب الطريق الكلي
        عبر $Q(3, 0)$ وتحقق من أنه يخسر أمام جوابك.
  \item صنّف الرباعي $A(1,1)$ و $B(4,2)$ و
        $C(5,5)$ و $D(2,4)$ تصنيفًا كاملًا: متوازي أضلاع؟ معيّن؟
        مستطيل؟ مربع؟
        (\cref{met:g10:coordgeom:parallelogram} و
        \cref{met:g10:coordgeom:triangle} --- وقارن القطرين
        أيضًا.)
  \item من أجل $A(-1, 2)$ و $B(3, 4)$ و $C(6, 0)$، أوجد $D$ بحيث
        يكون $ABCD$ متوازي أضلاع (حيلة \omterm{prop:g10:coordgeom:midpoint}{منتصف} القطر
        في مسألة نهاية الأسبوع عن أنصاف الدورات في الكتاب السابق، بالصيغ الآن).
  \item الخاتمة: ثلاث أدوات --- صيغة \omterm{prop:g10:coordgeom:midpoint}{المنتصف}، وصيغة
        المسافة، وطرح معادلتي دائرتين. سمِّ لكل منها
        نوع السؤال الذي حسمته في هذه المسألة، وبيّن
        ما يضيفه الفصلان التاليان إلى العدة
        (الاتجاهات والميول: المتجهات ومعادلات المستقيمات).
\end{enumerate}
\end{problem}
